\documentclass{exam}
\usepackage{../commonheader}

\discnumber{3}
\title{Higher-Order Functions \& Environment Diagrams}
\date{September 7, 2026 -- September 11, 2026}

\begin{document}
\maketitle

\begin{meta}
    \textbf{Example Timeline}
    \begin{itemize}
        \item Intros, Ice Breakers, Expectations, Logistics [5 min]
        \subitem Don't be afraid to spend some time on this! It'll likely make your job in the future easier! Prospective mentees will also get to know how CSM works and doing icebreakers will make it a welcoming experience for all of them but try to keep them kind of short since this is not their official session.
        \subitem Also, something we want to do this semester more of is emphasize the conceptual topics and concepts in 61A like functional abstraction and more. So make sure to talk about them regularly throughout all of your teaching sessions and tell junior mentors to also emphasize them as well.
        \item Environment Diagrams Mini Lecture + Q1 [10 mins]
        \subitem 1. You can use Q1 as an example for the mini-lecture
        \subitem 2. A lot of students (even those with significant programming background) get confused by env. diagrams. It's probably worth doing the minilecture even if you have advanced students.
        \subitem 3. You should address when we need to make a new frame in an environment diagram.
        \item Higher Order Functions Overview + Reasoning + Q1 + Q2 [15 mins]
        \item Problems (You pick which ones): HOF  
        \subitem 1. Foo: has an example of a function which is returned along with return in the middle of function and returning None
        \subitem 2. Compose: easy example to build skill in writing lambda functions.
        \subitem 3. WholeSum: Another HOF example and digit manipulation and while loop practice
        \subitem 4. Alternator: HOF example with while loop practice
        \subitem 5. Curryforever: An example with writing code for a function which takes in a variable number of arguments; the number of arguments that will be taken is input earlier
    \end{itemize}
\end{meta} 

\section{Environment Diagrams}
\begin{questions}
    \subimport{../../topics/functions-and-expressions/medium/env-diagram/}{joke.tex}
    \begin{questionmeta}
        Suggested Time: 10 min; Difficulty: Medium
        \begin{itemize}
            \item Make sure that the students understand how Python looks for a value of a variable, from local (to parent(s)) to global.
            \item Make sure your students understand the difference between an intrinsic name and a bound name
            \subitem Intrinsic: For user defined functions, this intrinsic name is the name used in the \lstinline{def} statement 
            \subitem Bound: Names of variables that point to the function object. A function can have many bound names, and the bound names of a function can often change. 
        \end{itemize}    
      \end{questionmeta}
          
\end{questions}

\section{Higher-Order Functions}
\begin{questions}
    \subimport{../../topics/hof/easy/}{why.tex}
    \begin{questionmeta}
        Suggested Time: 5 mins
        
        Treat this as more of a guiding question to start discussion on HOFs, like talk about it and do not treat it like other CSM questions. Try to give examples when explaining HOFs. Some examples are linked at this link: \url{https://tinyurl.com/5fsd3h8b}.
        One specific example to illustrate the use of lambdas is the optional \lstinline{key} parameter for \lstinline{min} and \lstinline{max} functions, but just as a side note you could also pass in a function defined using the keyword def for the \lstinline{key} parameter.
        Might want to show examples that illustrate each of the bullet points in action. 
    \end{questionmeta}
    
    \subimport{../../topics/hof/easy/env-diagram/}{foo.tex}
    \begin{questionmeta}
        Suggested Time: 7 mins; Difficulty: Easy
        \begin{itemize}
            \item Emphasize that if line with return statement is reached, then expression after return keyword will be evaluated and that value is then returned to the environment/frame in which the function returning the value was called in. 
            \item Also, emphasize that if no return statement is encountered while calling a function with certain arguments then the return value is None.
        \end{itemize}
    \end{questionmeta}

    \subimport{../../topics/environments/medium/}{abde.tex}
    \begin{questionmeta}
        Suggested Time: 10 min; Difficulty: Medium
        \begin{itemize}
            \item Make sure that students understand how lambda functions are expressed in the environment diagram.
            \item Make sure to stress variables in different scopes.
        \end{itemize}
      \end{questionmeta}

    \subimport{../../topics/hof/easy/}{compose.tex}
    \begin{questionmeta}
        Suggested Time: 5 mins; Difficulty: Easy
    \end{questionmeta}

    \subimport{../../topics/hof/medium/}{check_sum.tex}
    \begin{questionmeta}
        Suggested Time: 8 Mins; Difficulty: Medium
        \begin{itemize}
            \item Remind your students that for HOFs, you must \lstinline{return} the inner function (ie we must \lstinline{return check} to use it).
            \item Also depending on the skill level of students in your section, a recap of digit manipulation may be needed (ie x // 10, x \% 10, etc.)
        \end{itemize}
    \end{questionmeta}

    \subimport{../../topics/hof/medium}{alternator.tex}
    \begin{questionmeta}
        Suggested Time: 8 mins; Difficulty: Medium
        \begin{itemize}
            \item Walk students through each iteration from 1 to x, and show how each of the two functions \verb|f,g| alternate on incrementing inputs.
            \item Remember the general structure needed whenever a function must return a function.
        \end{itemize}
    \end{questionmeta}

    \subimport{../../topics/hof/medium/}{curry_forever.tex}
    \begin{questionmeta}
        Suggested Time: 15 mins; Difficulty: Medium
        \begin{itemize}
        \item Students might be confused on how to build a function which takes in a variable number of arguments. If they are stuck, give them some hints on how to approach in general questions with skeleton code.
        \end{itemize}
    \end{questionmeta}
\end{questions}

\end{document}